\section{A Frey Q-curve and a residual modular-level budget}
\label{fm-section}

This section gives a modular entry for the second mixed norm. Let
\[
 F=F(a,b)=a^4+3a^3b+5a^2b^2+3ab^3+b^4,
 \qquad a,b\in\mathbb Z_{>0},\quad\gcd(a,b)=1,
\]
and put $x=a^2+b^2$, $y=a+b$. We derive an exact conductor over
$\mathbb Q(\sqrt{-3})$ and a finite set of necessary modular levels for
the pure-power branch. The remaining newforms are not eliminated.
All representation-theoretic inputs are ordinary mathematical inputs;
this is not a full Lean formalization of Tate's algorithm or modularity.

\subsection{The actual generalized-Fermat coordinates}

\begin{lemma}[The square gate and two parity charts]\label{fm-square-gate}
One has
\begin{equation}\label{fm-coordinates}
 x^2+3y^4=4F,\qquad 2x-y^2=(a-b)^2,\qquad\gcd(x,y)\mid2.
\end{equation}
If $a,b$ have opposite parity, then $x,y,d=a-b$ are odd,
$\gcd(x,y)=1$, $d^2=2x-y^2$, and $|d|<y$. Conversely, these conditions
with $y>0$ recover a primitive positive seed by
$a=(y+d)/2$, $b=(y-d)/2$.

If $a,b$ are both odd, then
$X=x/2$, $Y=y/2$, $D=(a-b)/2$ satisfy
\begin{equation}\label{fm-odd-chart}
 X^2+12Y^4=F,\quad X-Y^2=D^2,\quad |D|<Y,
 \quad X\text{ odd},\quad\gcd(X,Y)=1.
\end{equation}
Conversely these conditions with $Y>0$ recover the primitive positive
odd seed $a=Y+D$, $b=Y-D$.
\end{lemma}
\begin{proof}
The identities follow by expansion. A common divisor of $x,y$ divides
$2ab$, while $\gcd(ab,a+b)=1$. In the first inverse, an odd common
prime of $a,b$ would divide $x,y$, and both seeds cannot be even since
$x$ is odd. In the second inverse, $X=Y^2+D^2$ odd makes $Y,D$
opposite in parity, and $\gcd(X,Y)=1$ implies $\gcd(Y,D)=1$.
The strict inequalities give positivity in both cases.
\end{proof}

An arbitrary solution of $x^2+3y^4=4VQ^g$ is therefore not yet an
actual seed: it must satisfy the additional square and positivity gate.
Also, the known equation $A^4+3B^2=C^p$ has different coefficients.
The substitution $A=3y,B=3x$ gives $A^4+3B^2=108VQ^g$ with non-coprime
parameters, so its equation-specific non-existence results do not apply.

The quartic $F$ is odd and $F\equiv1\pmod3$. The latter follows from
$F\equiv(a^2+b^2)^2\pmod3$; in particular $3\nmid x$.
Lemma~\ref{qc-bad-depth} gives the exact additional fact
$13\mid F\Rightarrow v_{13}(F)=1$. Thus
\begin{equation}\label{fm-pure-support}
 F=Q^g,\quad g\ge2\quad\Longrightarrow\quad\gcd(Q,78)=1.
\end{equation}

\subsection{An explicit degree-two Q-curve}

Put $r=\sqrt{-3}$, $K=\mathbb Q(r)$ and
$\alpha_\pm=3y^2\pm xr$. Attach
\begin{equation}\label{fm-frey-curve}
 E_{a,b}:\quad \mathcal Y^2=\mathcal X^3+
                 12y\mathcal X^2+6\alpha_+\mathcal X.
\end{equation}

\begin{lemma}[Isogeny and invariant identities]\label{fm-invariants}
The curve~\eqref{fm-frey-curve} is a degree-two Q-curve, with the
displayed isogenies defined over $K(\sqrt{-2})$. Its invariants are
\begin{equation}\label{fm-invariant-formulas}
 \begin{aligned}
 \alpha_+\alpha_-&=12F,&
 c_4&=2^5 3^2(5y^2-xr),\\
 c_6&=2^8 3^3y(3xr-7y^2),&
 \Delta&=2^{11}3^4F\alpha_+.
 \end{aligned}
\end{equation}
At every rational prime $q>3$ not dividing $F$ it has good reduction.
If $q\mid F$, then $q$ splits in $K$, the displayed equation is minimal
and multiplicative at both primes, and, for $m=v_q(F)$,
\begin{equation}\label{fm-split-valuations}
 \{v_{\mathfrak q}(\Delta_{\min}),
       v_{\bar{\mathfrak q}}(\Delta_{\min})\}=\{m,2m\}.
\end{equation}
\end{lemma}
\begin{proof}
Write $A=12y$ and $B=6\alpha_+$. The quotient by the order-two point
$(0,0)$ has coefficients $-2A=-24y$ and
$A^2-4B=24\alpha_-$. Hence it is the quadratic twist by $-2$ of the
conjugate curve. This is the explicit isogeny construction also used in
\cite[Proposition~2.2]{PacettiVillagra2022QCurves}; it does not require a
primitive solution of the equation studied there.
The formulas
$c_4=16(A^2-3B)$, $c_6=-64A^3+288AB$ and
$\Delta=16B^2(A^2-4B)$ prove~\eqref{fm-invariant-formulas}.

If $q\nmid F$, the discriminant is a unit. If $q\mid F$, neither $x$
nor $y$ is divisible by $q$: otherwise~\eqref{fm-coordinates} would
force both, contradicting $\gcd(x,y)\mid2$. The same equation makes
$-3$ a nonzero square modulo $q$, so $q$ splits. The elements
$\alpha_+,\alpha_-$ cannot both vanish at a prime above $q$, since their
sum and difference would force $x,y$ to vanish. Their local valuations
are $m,0$ in one order. The factor $5y^2-xr$ in $c_4$ reduces to
$8y^2$ or $2y^2$, a unit. Thus the equation is minimal and multiplicative,
with the two valuations in~\eqref{fm-split-valuations}.

Since $F\ge13$ and $\gcd(F,6)=1$, such a multiplicative prime always
exists. Its negative valuation of $j$ excludes complex multiplication,
because CM $j$-invariants are algebraic integers. In the non-CM case the
isogenies to the conjugate form a rank-one module, so their degrees are
the minimum degree times integer squares. The existing degree-two
isogeny forces the minimum degree to be exactly two.
\end{proof}

\subsection{Uniform Tate calculations at two and three}

Let $\mathfrak p_2=(2)$, an inert prime with residue field $\mathbb F_4$,
and $\mathfrak p_3=(r)$, with $v_{\mathfrak p_3}(3)=2$.
We use valuations normalized by $v_{\mathfrak p_2}(2)=1$ and
$v_{\mathfrak p_3}(r)=1$.

\begin{theorem}[The exact conductor over the quadratic field]
\label{fm-quadratic-conductor}
The curve has minimal discriminant valuation $12$, Kodaira type $I_2^*$
and conductor exponent $6$ at $\mathfrak p_2$. At $\mathfrak p_3$ the
corresponding data are $9$, $III^*$ and $2$. Consequently
\begin{equation}\label{fm-exact-conductor}
 \mathcal N(E)=\mathfrak p_2^6\mathfrak p_3^2
       \prod_{q\mid F,\ q>3}\mathfrak q\bar{\mathfrak q}.
\end{equation}
\end{theorem}
\begin{proof}
At $\mathfrak p_3$, $3\nmid x$ gives $v(\alpha_+)=1$, hence
$v(\Delta)=9$. The coefficients satisfy
\[
 a_1=a_3=a_6=0,\qquad v(a_2)\ge2,\qquad v(a_4)=3.
\]
The auxiliary cubic in Tate's additive normalization is $T^3$.
After its zero triple root, the quadratic for the $IV^*$ test is $T^2$.
The stronger divisibility of $a_3,a_6$ after its zero repeated root
already holds. The next test is nonzero since $a_4/r^3$ is a unit;
thus the type is $III^*$. Its eight geometric components give conductor
exponent $9+1-8=2$. These branch tests also establish minimality.

At $\mathfrak p_2$, put $\zeta=(1+r)/2$. For integers $U,V$ and $n\ge1$,
the integral-basis identity $U+Vr=(U-V)+2V\zeta$ gives
\begin{equation}\label{fm-integral-basis-test}
 U+Vr\in2^n\mathcal O_{K,\mathfrak p_2}
 \ \Longleftrightarrow\ 2^n\mid U-V,\quad 2^{n-1}\mid V.
\end{equation}
Conjugation preserves the valuation and $\alpha_+\alpha_-=12F$, so
$v(\alpha_+)=1$ and $v(\Delta)=12$.
Make the integral substitution $\mathcal X=X'+R$, $\mathcal Y=Y'+T$
specified by the following table.
\[
\begin{array}{c|c|c}
\text{condition}&R&T\\ \hline
y\equiv0\pmod4&2&2(r+1)\\
y\equiv2\pmod4&2&2(r-1)\\
y\text{ odd, even seed }0\pmod4&1+r&4\\
y\text{ odd, even seed }2\pmod4&1+r&2(r+1)
\end{array}
\]
The transformed coefficients are
\begin{equation}\label{fm-transformed-coefficients}
\begin{aligned}
 a'_1&=0,&a'_2&=A+3R,&a'_3&=2T,\\
 a'_4&=B+2AR+3R^2,&
 a'_6&=R^3+AR^2+BR-T^2.
\end{aligned}
\end{equation}
We verify the uniform pattern
\begin{equation}\label{fm-tate-pattern}
 v(a'_2)=1,\quad v(a'_3)\ge3,\quad v(a'_4)=3,\quad v(a'_6)\ge5.
\end{equation}

First suppose $y$ is even. Then $x\equiv2\pmod8$ and $a'_2=12y+6$
has valuation one. With $T=2(r+s)$, where $s=1$ or $-1$ according to
the table, write
\[
 a'_4=6(H+xr),\quad H=3y^2+8y+2,
 \qquad
 a'_6=C+Dr,
\]
where $C=16+48y+36y^2$ and $D=12x-8s$.
Both $H/2,x/2$ are odd. For odd integers $u,v$, the norm
$u^2+3v^2\equiv4\pmod8$ makes $v_{\mathfrak p_2}(u+vr)=1$.
This proves $v(a'_4)=3$. Also $16\mid D$ and
\[
 (C-D)/4=4+12y+9y^2-3x+2s\equiv0\pmod8,
\]
using $x\equiv2\pmod8$ and the two even values of $y\pmod4$.
Equation~\eqref{fm-integral-basis-test} gives $v(a'_6)\ge5$.
Finally $v(T)=2$, giving $v(a'_3)=3$.

If $y$ is odd, then
\[
 a'_2=3(4y+1+r),\quad a'_4=6(H+Jr),\quad
 H=3y^2+4y-1,\quad J=x+4y+1.
\]
The first bracket has valuation one; $H/2,J/2$ are odd, so again
$v(a'_4)=3$. For $T=4$, the coefficients of $a'_6=C+Dr$ obey
\[
 D=24y+18y^2+6x,\qquad C-D=-24(1+2y+x).
\]
Here $x\equiv1\pmod8$ makes $16\mid D$, and $x\equiv1\pmod4$ gives
$32\mid C-D$. For $T=2(r+1)$ the corresponding formulas are
\[
 D=24y+18y^2+6x-8,\qquad C-D=8(1-6y-3x).
\]
Now $x\equiv5\pmod8$ proves $16\mid D$, and again $32\mid C-D$.
In both cases $v(a'_3)=3$, completing~\eqref{fm-tate-pattern}.

Tate's auxiliary cubic is $Z^2(Z+b)$, where $b=a'_2/2$ is nonzero in
$\mathbb F_4$. Its double root is zero. The first quadratic in the
$I_n^*$ branch is $Z^2$, since $a'_3/4,a'_6/16$ vanish in the residue
field. No further translation is needed: $8\mid a'_3$ and $32\mid a'_6$
already hold. The second quadratic is
\[
 (a'_2/2)Z^2+(a'_4/8)Z+a'_6/32,
\]
whose derivative $a'_4/8$ is a unit. The algorithm stops at $I_2^*$;
its seven geometric components give conductor exponent $12+1-7=6$.
All changes have scaling one. These are residue-field polynomial tests,
valid over $\mathbb F_4$; no shortcut peculiar to $\mathbb F_2$ is used.
For the explicit double-root and triple-root branches see
\cite[Chapter~3, pp.~66--68]{Cremona1997Algorithms}.
\end{proof}

\subsection{Five necessary levels in the pure-power branch}

\begin{theorem}[Residual modular levels]\label{fm-five-levels}
If $F(a,b)=Q^p$ for $Q\in\mathbb Z_{>0}$ and a prime $p>7$, then
there is a weight-two newform
with character $\epsilon=\psi_3$ of conductor $12$, and level
\begin{equation}\label{fm-five-level-list}
 N=2^e3^2,\quad2\le e\le6,
 \qquad N\in\{36,72,144,288,576\},
\end{equation}
whose residual representation restricts to the specified twist of
$E_{a,b}[p]$ over $K$.
\end{theorem}
\begin{proof}
We spell out the external representation inputs because the level formula
for the different Fermat equation cannot be copied.
By \cite[Theorem~3.2 and Table~3.3]{PacettiVillagra2022QCurves}, with
$d=3$, there is a fixed order-four character $\chi$ of $G_K$ satisfying
\[
 \operatorname{cond}(\chi)\mid\mathfrak p_2^3,\qquad
 \chi^2=\epsilon|_{G_K},\qquad
 {}^\sigma\chi=\chi\psi_{-2}|_{G_K}.
\]
In particular it is unramified at three. This character construction
does not require a primitive Fermat solution. The isogeny of
Lemma~\ref{fm-invariants} makes $\rho_E\otimes\chi$ invariant under
quadratic conjugation. The extension proof in
\cite[Theorem~4.2]{PacettiVillagra2022QCurves} applies to this non-CM
Q-curve. Its oddness invokes the general Q-curve result of Ribet cited
there; it does not follow from Schur's lemma alone. The determinant is
$\epsilon\chi_{\mathrm{cyc}}$.

At every $q>3$, Lemma~\ref{fm-invariants} makes the minimal discriminant
valuations divisible by $p$. The Tate-curve criterion gives unramified
$p$-torsion away from $6p$ and finite flat $p$-torsion at $p$.
Indeed, at $p$ the Tate parameter has valuation divisible by $p$, so
its Kummer class has a unit representative; good reduction gives finite
flatness directly. The extension $K/\mathbb Q$ and the character $\chi$
are unramified at $p$, preserving these properties under twist and descent.
The same local mechanism is explained in
\cite[Propositions~3.6--3.9]{Koutsianas2020GeneralizedFermat}.

At $\mathfrak p_2$, the curve is potentially good: the invariant formulas
give $v(c_4)=6$, $v(\Delta)=12$, hence $v(j)=6>0$.
On every finite higher inertia group its characteristic-zero
representation has determinant one. Its fixed subspace cannot have
dimension one: semisimplicity would make the complementary character
equal to the determinant, hence trivial as well. Thus its ramification
has a single break. The conductor exponent $6$ gives break $2$.
The character $\chi$ has break at most $2$, so the twist has conductor
exponent at most $6$. Since $K_{\mathfrak p_2}/\mathbb Q_2$ is unramified,
restriction preserves that exponent. Its determinant has conductor
exponent $2$, furnishing the lower bound $e\ge2$.

At three, the curve is potentially good and has conductor exponent two,
so it has no wild ramification. Twisting by $\chi$ preserves this.
The quadratic extension $K_{\mathfrak p_3}/\mathbb Q_3$ is tame, so both
extensions to $G_{\mathbb Q_3}$ have trivial wild inertia and conductor
exponent at most two. They differ by $\psi_{-3}$. The conductor of their
direct sum, the induced representation, has exponent
$2\cdot1+1\cdot2=4$ by
\cite[Theorem~4.1]{PacettiVillagra2022QCurves}. Hence each exponent
equals two.

The big-image input is the general Q-curve statement
\cite[Theorem~5.2]{PacettiVillagra2022QCurves}: squarefree isogeny degree
and a multiplicative prime greater than three suffice. Lemma~\ref{fm-invariants}
checks both conditions, independently of the nonprimitive substitution
$(3y,3x)$. For $K=\mathbb Q(\sqrt{-3})$ the source gives $N_3=7$;
the improvement uses
\cite[Proposition~5.4]{Koutsianas2020GeneralizedFermat}.
This is a published external input, not a new modular-symbol certificate
in this repository.

Large image, oddness and finite flatness give Serre weight two and the
listed prime-to-$p$ levels, by the modularity and lowering arguments in
the cited sources. The finite inertial images at two and three have
order supported at two and three; adjoining the order-four character
and the quadratic descent does not change that support. Reduction
modulo $p>7$ therefore preserves their conductors and the character
$\epsilon$. This proves the necessary level list.
\end{proof}

\subsection{A direct budget for moving residual support}

\begin{theorem}[Residual level and dimension budgets]\label{fm-residual-budget}
Suppose $F=VQ^p$, where $V,Q\in\mathbb Z_{>0}$, $p>7$ is prime,
$V$ is $p$-power-free, and
$p\nmid V$. Then the necessary weight-two newform has character
$\epsilon=\psi_3$ and exact prime-to-$p$ level
\begin{equation}\label{fm-residual-level}
 N=2^e3^2\operatorname{rad}(V),\qquad2\le e\le6.
\end{equation}
For $\lambda=\max\{1,\log V\}/p$ this implies
\begin{equation}\label{fm-log-level}
 \frac{\log N}{p}\le\lambda+\frac{\log576}{p}.
\end{equation}
For all $1\le V\le\exp(Lp)$ together, with $L\ge0$, the sum of the
dimensions of the five possible newspaces is at most
$965\exp(3Lp)$. In particular this budget is subexponential when
$L\to0$ and $p\to\infty$.
\end{theorem}
\begin{proof}
At $q>3$, write $m=v_q(V)+p v_q(Q)$. The two local minimal discriminant
valuations are $m,2m$. The Tate residual inertia is nontrivial precisely
when $p$ does not divide its parameter's valuation. Since $p$ is odd and
$V$ is $p$-power-free, this occurs precisely at $q\mid V$.
The twist is unramified at these primes. At $p$, the condition $p\nmid V$
ensures finite flatness. All other steps in Theorem~\ref{fm-five-levels}
are unchanged, proving~\eqref{fm-residual-level}.
Then $N\le576V$ gives~\eqref{fm-log-level}.

Let $D=\operatorname{rad}(V)$; it is prime to six. The index satisfies
\[
 I=[\mathrm{SL}_2(\mathbb Z):\Gamma_0(N)]
  =2N\prod_{q\mid D}(1+q^{-1})
  \le2ND\le1152V^2.
\]
The weight-two Sturm bound gives dimension at most
$\lfloor I/6\rfloor+1\le193V^2$ for each space. The five dimensions
sum to at most $965V^2$. Put $H=\lfloor\exp(Lp)\rfloor$.
Even allowing inadmissible $V$, there are at most $H$ terms of size at
most $965H^2$, giving the asserted bound. Dimensions count a Galois
orbit with its coefficient-field degree, not merely as one orbit.
The index and Sturm-bound conventions agree with
\cite[\texttt{mfsturm}]{PARIModularForms}.
\end{proof}

The condition $p\nmid V$ cannot be inferred from a small $\lambda$;
Section~\ref{fm7-section} includes the non-finite-flat branch $p\mid V$
at weight $p+1$, with uniform candidate elimination still open.
The level budget does not make the newspaces empty and gives no seed
height bound. Even in the pure branch the five spaces must be examined
against the actual seed family and its inverse square gate.
If a prime $p$ divides a larger pure exponent $g$, write
$Q^g=(Q^{g/p})^p$. An eventual contradiction for all large $p$, together
with the earlier fixed-exponent finiteness, would give finiteness for
pure-power seeds admitting an exponent $g\ge3$. No such contradiction
is proved here. The known infinite square family is outside that
proposed implication. The moving Kummer-cover and point-height routes
remain open alongside this modular entry.
